\documentclass[11pt]{article}
\usepackage{reusv}
\usepackage{listings}
\usepackage{enumitem}

\lstset{basicstyle=\ttfamily\small, frame=single, breaklines=true}
\setborderthickness{0.5pt}
\setbordermargin{1cm}
\setbordercolor{black}

\slimtitle{A3. 호만 전이 해석해 및 검증 벤치마크}

\begin{document}
\drawborder
\thispagestyle{firstpage}

\lightertitle{A3. 호만 전이 해석해 및 검증 벤치마크}
\def\lighterdate{February 12, 2026}
\lighterauthor{Suwon Lee}
\lighteraddress{}{Future Mobility Control Lab, Department of Future Mobility, Kookmin University, Seoul, Republic of Korea}
\lighteremail{suwon@kookmin.ac.kr}

\begin{abstract}
호만 전이(Hohmann transfer)는 두 동심 원궤도 사이의 이-임펄스(two-impulse) 최소 $\Delta v$ 궤도전이이다.
본 보고서는 vis-viva 방정식으로부터 호만 전이 공식을 유도하고, 코드 구현과의 일대일 대응을 보이며, 수치 검증 결과를 제시한다.
호만 전이 해석해는 연속 추력 최적화 결과의 물리적 타당성 검증 기준(하한)과 궤적 최적화 코드의 단위 테스트 기준값으로 활용된다.
\end{abstract}

% ============================================================================
\section{목적}
% ============================================================================

호만 전이(Hohmann transfer)는 두 동심 원궤도 사이의 이-임펄스(two-impulse) 최소 $\Delta v$ 궤도전이이다. 본 프로젝트에서 호만 전이 해석해는 두 가지 역할을 수행한다.

첫째, 연속 추력 최적화 결과의 \textbf{물리적 타당성 검증 기준}이다. 연속 추력 해의 총 $\Delta v$는 이상적 임펄스 해의 $\Delta v$보다 항상 크거나 같아야 하므로, 호만 전이 $\Delta v$는 하한(lower bound)을 제공한다.

둘째, 궤적 최적화 코드의 \textbf{단위 테스트 기준값}이다. LEO-GEO 호만 전이처럼 교과서에서 널리 알려진 참조값과 비교하여 구현 정확도를 검증한다.

본 보고서는 vis-viva 방정식으로부터 호만 전이 공식을 유도하고, 코드 구현과의 일대일 대응을 보이며, 수치 검증 결과를 제시한다.

\vspace{0.5em}
\noindent 대응 코드: \texttt{src/orbit\_transfer/astrodynamics/hohmann.py}

% ============================================================================
\section{수학적 배경}
% ============================================================================

\subsection{Vis-viva 방정식}

이체 문제(two-body problem)에서 우주체의 비역학 에너지(specific mechanical energy)는 궤도상 어디서나 보존된다. 위치 $r$, 속력 $v$에서
\[
\varepsilon = \frac{v^2}{2} - \frac{\mu}{r}
\]
이며, 이 값은 궤도 장반경 $a$로만 결정된다:
\[
\varepsilon = -\frac{\mu}{2a}
\]

두 식을 연립하면 \textbf{vis-viva 방정식}을 얻는다:
\[
v^2 = \mu\left(\frac{2}{r} - \frac{1}{a}\right)
\]

이 관계는 케플러 궤도상 임의 위치에서의 속력을 궤도 크기($a$)와 현재 위치($r$)만으로 결정한다. 호만 전이의 모든 속도 계산은 이 단일 방정식에 기초한다~\cite{Curtis2020}.

\subsection{원궤도 속도}

원궤도에서는 $r = a$이므로 vis-viva 방정식이 단순화된다:
\[
v_{\text{circ}} = \sqrt{\frac{\mu}{a}}
\]

초기 원궤도(반지름 $r_1 = a_1$)와 최종 원궤도(반지름 $r_2 = a_2$)의 속도는 각각
\[
v_{c1} = \sqrt{\frac{\mu}{a_1}}, \qquad v_{c2} = \sqrt{\frac{\mu}{a_2}}
\]
이다.

\subsection{호만 전이 궤도 설계}

호만 전이 궤도는 근지점이 초기 궤도에, 원지점이 최종 궤도에 접하는 타원이다~\cite{Hohmann1960}. 따라서 전이 타원의 장반경은
\[
a_t = \frac{r_1 + r_2}{2} = \frac{a_1 + a_2}{2}
\]
이다. 전이 궤도상의 속도를 vis-viva 방정식으로 구한다.

\textbf{근지점 속도} ($r = r_1 = a_1$):
\[
v_{t1} = \sqrt{\mu\left(\frac{2}{a_1} - \frac{1}{a_t}\right)}
\]

\textbf{원지점 속도} ($r = r_2 = a_2$):
\[
v_{t2} = \sqrt{\mu\left(\frac{2}{a_2} - \frac{1}{a_t}\right)}
\]

\subsection{$\Delta v$ 공식}

호만 전이는 두 번의 접선 임펄스 기동으로 구성된다.

\textbf{제1 기동 (근지점, 초기 궤도 $\rightarrow$ 전이 궤도):}
\[
\Delta v_1 = v_{t1} - v_{c1} = \sqrt{\mu\left(\frac{2}{a_1} - \frac{1}{a_t}\right)} - \sqrt{\frac{\mu}{a_1}}
\]

궤도를 올리는 경우($a_2 > a_1$) $v_{t1} > v_{c1}$이므로 $\Delta v_1 > 0$이다. 순방향 접선 가속이다.

\textbf{제2 기동 (원지점, 전이 궤도 $\rightarrow$ 최종 궤도):}
\[
\Delta v_2 = v_{c2} - v_{t2} = \sqrt{\frac{\mu}{a_2}} - \sqrt{\mu\left(\frac{2}{a_2} - \frac{1}{a_t}\right)}
\]

원지점에서 전이 궤도 속도 $v_{t2}$는 최종 원궤도 속도 $v_{c2}$보다 작으므로 $\Delta v_2 > 0$이다.

\textbf{총 $\Delta v$:}
\[
\Delta v_{\text{total}} = |\Delta v_1| + |\Delta v_2|
\]

구현에서는 절대값을 사용하여 궤도를 내리는 경우($a_2 < a_1$)도 동일한 함수로 처리한다. 이 경우 $\Delta v_1 < 0$ (감속), $\Delta v_2 < 0$ (감속)이나 기동 크기의 합은 동일하다.

\subsection{비행시간}

전이 궤도의 반주기가 곧 비행시간이다. 케플러 제3법칙에 의해
\[
T = 2\pi\sqrt{\frac{a_t^3}{\mu}}
\]
이므로 반주기(근지점 $\rightarrow$ 원지점)는
\[
\text{TOF} = \pi\sqrt{\frac{a_t^3}{\mu}}
\]
이다~\cite{Curtis2020}.

\subsection{접선 기동의 최적성과 Oberth 효과}

호만 전이가 이-임펄스 전이 중 최소 $\Delta v$를 달성하는 이유는 다음 두 원리에 기인한다.

\textbf{접선 기동의 에너지 효율.} 임펄스 $\Delta \mathbf{v}$가 속도 벡터에 평행하게 가해질 때 궤도 에너지 변화가 최대이다. 에너지 변화율은
\[
\Delta \varepsilon = \mathbf{v} \cdot \Delta \mathbf{v} + \frac{|\Delta \mathbf{v}|^2}{2}
\]
이므로, 동일한 $|\Delta \mathbf{v}|$에 대해 $\mathbf{v} \parallel \Delta \mathbf{v}$일 때 $\Delta \varepsilon$가 최대화된다~\cite{Lawden1963}.

\textbf{Oberth 효과.} 속력이 큰 위치에서 동일한 $|\Delta v|$를 가하면 운동에너지 변화($\Delta E_k = mv\Delta v + \frac{1}{2}m\Delta v^2$)가 더 크다. 호만 전이는 제1 기동을 근지점(속력 최대)에서 수행하여 이 효과를 활용한다. Edelbaum~\cite{Edelbaum1967}은 동심 원궤도 전이에서 이-임펄스 호만 전이가 최적임을 증명하였으며, 궤도 반지름 비 $a_2/a_1$이 약 11.94를 초과하면 삼-임펄스 bi-elliptic 전이가 더 효율적임을 보였다.

\subsection{연속 추력 $L_2$ 비용과 임펄스 $\Delta v$의 관계}

본 프로젝트의 궤적 최적화에서 사용하는 비용함수는 제어 입력의 $L_2$ 노름 제곱 적분이다:
\[
J = \int_0^{t_f} \|\mathbf{u}(t)\|^2 \, dt
\]

여기서 $\mathbf{u}(t)$는 추력 가속도 벡터이다. 임펄스 기동은 Dirac delta 형태의 제어 입력으로 표현된다:
\[
\mathbf{u}(t) = \sum_{k} \Delta \mathbf{v}_k \, \delta(t - t_k)
\]

임펄스 해는 연속 추력 해의 극한 형태로 해석할 수 있다. 연속 추력 해의 해공간 $\mathcal{S}_{\text{cont}}$는 임펄스 해의 해공간 $\mathcal{S}_{\text{imp}}$를 포함한다:
\[
\mathcal{S}_{\text{imp}} \subset \mathcal{S}_{\text{cont}}
\]

이 포함관계는 연속 추력 최적해의 비용이 임펄스 최적해의 비용보다 항상 작거나 같음을 의미한다~\cite{Taheri2018, Prussing2010}. 따라서 호만 전이 $\Delta v$는 연속 추력 해의 $\Delta v$ 하한을 제공하며, 이를 위반하는 수치 해는 물리적으로 부당하다.

구체적으로, 추력 크기 $T_{\max}$와 비행시간 $t_f$가 충분히 클 때, 연속 추력 최적 프로파일은 기동 구간이 집중된 ``bang-off-bang'' 형태로 수렴하며, 이는 임펄스 해에 접근한다~\cite{Prussing2010}. 본 프로젝트에서는 이 수렴 거동을 데이터베이스 전체에서 체계적으로 확인하는 데 호만 전이 해석해를 활용한다.

% ============================================================================
\section{구현 매핑}
% ============================================================================

\texttt{hohmann.py}는 두 개의 함수로 구성되며, 각 함수는 2.3--2.5절의 수식을 직접 구현한다. 사용 상수는 \texttt{constants.py}에 정의된 $\mu_\oplus = 398600.4418$~km$^3$/s$^2$, $R_\oplus = 6378.137$~km이다.

\subsection{함수별 매핑}

\begin{table}[htbp]
\centering
\small
\begin{tabular}{@{}llll@{}}
\toprule
수학 표현 & Python 구현 & 파일 위치 & 비고 \\
\midrule
$a_t = (a_1 + a_2)/2$ & \texttt{a\_t = (a1 + a2) / 2.0} & hohmann.py:28 & 전이 타원 장반경 \\
$v_{c1} = \sqrt{\mu/a_1}$ & \texttt{v1 = np.sqrt(mu / a1)} & hohmann.py:31 & 초기 원궤도 속도 \\
$v_{c2} = \sqrt{\mu/a_2}$ & \texttt{v2 = np.sqrt(mu / a2)} & hohmann.py:32 & 최종 원궤도 속도 \\
$v_{t1} = \sqrt{\mu(2/a_1 - 1/a_t)}$ & \texttt{v\_t1 = np.sqrt(mu * (2.0/a1 - 1.0/a\_t))} & hohmann.py:35 & 전이 근지점 속도 \\
$v_{t2} = \sqrt{\mu(2/a_2 - 1/a_t)}$ & \texttt{v\_t2 = np.sqrt(mu * (2.0/a2 - 1.0/a\_t))} & hohmann.py:36 & 전이 원지점 속도 \\
$\Delta v_1 = |v_{t1} - v_{c1}|$ & \texttt{dv1 = abs(v\_t1 - v1)} & hohmann.py:39 & 절대값 처리 \\
$\Delta v_2 = |v_{c2} - v_{t2}|$ & \texttt{dv2 = abs(v2 - v\_t2)} & hohmann.py:40 & 절대값 처리 \\
$\Delta v_{\text{total}} = \Delta v_1 + \Delta v_2$ & \texttt{dv\_total = dv1 + dv2} & hohmann.py:41 & 총 기동량 \\
$\text{TOF} = \pi\sqrt{a_t^3/\mu}$ & \texttt{tof = np.pi * np.sqrt(a\_t**3 / mu)} & hohmann.py:64 & 전이 반주기 \\
\bottomrule
\end{tabular}
\end{table}

\subsection{함수 인터페이스}

\textbf{\texttt{hohmann\_dv(a1, a2, mu)}}

\begin{itemize}[nosep]
\item 입력: 초기 원궤도 반지름 $a_1$~[km], 최종 원궤도 반지름 $a_2$~[km], 중력 파라미터 $\mu$~[km$^3$/s$^2$]
\item 출력: \texttt{(dv1, dv2, dv\_total)}~[km/s]
\item $a_1 = a_2$이면 $\Delta v_1 = \Delta v_2 = 0$을 반환한다.
\item $a_2 < a_1$ (궤도 하강)에서도 절대값 처리로 양수 $\Delta v$를 반환한다.
\end{itemize}

\textbf{\texttt{hohmann\_tof(a1, a2, mu)}}

\begin{itemize}[nosep]
\item 입력: 동일
\item 출력: 비행시간 $\text{TOF}$~[s]
\item 전이 타원의 반주기로 계산한다.
\end{itemize}

\subsection{설계 결정}

\begin{enumerate}[nosep]
\item \textbf{절대값 사용}: $\Delta v_1$, $\Delta v_2$ 계산에 \texttt{abs()}를 적용하여 궤도 상승/하강 모두 동일 함수로 처리한다. 이는 호만 전이의 대칭성($a_1 \leftrightarrow a_2$이면 $\Delta v_{\text{total}}$ 동일)을 보존한다.
\item \textbf{NumPy 사용}: \texttt{np.sqrt}를 사용하여 배열 입력에도 대응 가능하나, 현재는 스칼라 입력만 사용한다.
\item \textbf{$\mu$ 외부 전달}: 중력 파라미터를 함수 인자로 받아 지구 외 천체에도 적용 가능하다.
\end{enumerate}

% ============================================================================
\section{수치 검증}
% ============================================================================

\subsection{LEO $\rightarrow$ GEO 벤치마크}

가장 널리 알려진 호만 전이 벤치마크인 LEO-GEO 전이를 검증 기준으로 사용한다~\cite{Curtis2020}.

\textbf{입력 조건:}

\begin{table}[htbp]
\centering
\begin{tabular}{@{}lll@{}}
\toprule
파라미터 & 값 & 비고 \\
\midrule
초기 고도 $h_1$ & 200~km & LEO \\
최종 고도 $h_2$ & 35,786~km & GEO \\
$a_1 = R_\oplus + h_1$ & 6,578.137~km & \\
$a_2 = R_\oplus + h_2$ & 42,164.137~km & \\
$\mu$ & 398,600.4418~km$^3$/s$^2$ & \\
\bottomrule
\end{tabular}
\end{table}

\textbf{계산 과정:}

전이 타원 장반경:
\[
a_t = \frac{6578.137 + 42164.137}{2} = 24371.137 \text{~km}
\]

원궤도 속도:
\[
v_{c1} = \sqrt{\frac{398600.4418}{6578.137}} = 7.784 \text{~km/s}
\]
\[
v_{c2} = \sqrt{\frac{398600.4418}{42164.137}} = 3.075 \text{~km/s}
\]

전이 궤도 속도:
\[
v_{t1} = \sqrt{398600.4418 \times \left(\frac{2}{6578.137} - \frac{1}{24371.137}\right)} = 10.239 \text{~km/s}
\]
\[
v_{t2} = \sqrt{398600.4418 \times \left(\frac{2}{42164.137} - \frac{1}{24371.137}\right)} = 1.597 \text{~km/s}
\]

$\Delta v$ 결과:
\[
\Delta v_1 = 10.239 - 7.784 = 2.455 \text{~km/s}
\]
\[
\Delta v_2 = 3.075 - 1.597 = 1.478 \text{~km/s}
\]
\[
\Delta v_{\text{total}} = 2.455 + 1.478 \approx 3.935 \text{~km/s}
\]

비행시간:
\[
\text{TOF} = \pi\sqrt{\frac{24371.137^3}{398600.4418}} \approx 18{,}925 \text{~s} \approx 5.26 \text{~hours}
\]

\textbf{검증 기준:} 교과서 참조값 $\Delta v_{\text{total}} \approx 3.935$~km/s~\cite{Curtis2020}와 0.01~km/s 이내 일치.

\subsection{동일 궤도 전이 (퇴화 케이스)}

$a_1 = a_2$이면 $a_t = a_1$이므로 $v_{t1} = v_{c1}$, $v_{t2} = v_{c2}$이다. 따라서
\[
\Delta v_1 = \Delta v_2 = 0, \qquad \Delta v_{\text{total}} = 0
\]

이 케이스는 구현의 수치적 안정성(0 반환)을 확인한다.

\subsection{테스트 구현}

검증은 \texttt{tests/test\_orbital\_elements.py}의 \texttt{TestHohmann} 클래스에서 수행된다.

\textbf{\texttt{test\_leo\_to\_geo}}: LEO(200~km) $\rightarrow$ GEO(35,786~km) 전이에서 $|\Delta v_{\text{total}} - 3.935| < 0.01$~km/s를 확인하고, $\Delta v_1 > 0$, $\Delta v_2 > 0$을 검증한다.

\textbf{\texttt{test\_hohmann\_tof}}: 동일 조건에서 비행시간이 양수이며 약 5.0--5.5시간 범위에 있음을 확인한다.

\subsection{검증 결과 요약}

\begin{table}[htbp]
\centering
\begin{tabular}{@{}llll@{}}
\toprule
검증 항목 & 기대값 & 허용 오차 & 결과 \\
\midrule
$\Delta v_{\text{total}}$ (LEO$\rightarrow$GEO) & 3.935~km/s & $\pm$~0.01~km/s & 통과 \\
$\Delta v_1, \Delta v_2 > 0$ & 양수 & -- & 통과 \\
TOF (LEO$\rightarrow$GEO) & 약 5.26~hours & 5.0--5.5~hours & 통과 \\
$a_1 = a_2$ 시 $\Delta v = 0$ & 0 & 수치 영점 & 통과 \\
\bottomrule
\end{tabular}
\end{table}

% ============================================================================
\bibliographystyle{IEEEtran}
\bibliography{references}
\end{document}
